\begin{lemma}
\label{lemma-weak-dimension-at-most-1}
Let $A$ be a ring. The following are equivalent
\begin{enumerate}
\item $A$ has weak dimension $\leq 1$,
\item every ideal of $A$ is flat,
\item every finitely generated ideal of $A$ is flat,
\item every submodule of a flat $A$-module is flat, and
\item every local ring of $A$ is a valuation ring.
\end{enumerate}
\end{lemma}

\begin{proof}
If $A$ has weak dimension $\leq 1$, then the resolution
$0 \to I \to A \to A/I \to 0$ shows that every ideal $I$
is flat by Lemma \ref{lemma-last-one-flat}.
Hence (1) $\Rightarrow$ (2).

\medskip\noindent
Assume (4). Let $M$ be an $A$-module. Choose a surjection
$F \to M$ where $F$ is a free $A$-module. Then $\Ker(F \to M)$
is flat by assumption, and we see that $M$ has tor dimension
$\leq 1$ by Lemma \ref{lemma-tor-dimension}.
Hence (4) $\Rightarrow$ (1).

\medskip\noindent
Every ideal is the union of the finitely generated ideals
contained in it. Hence (3) implies (2) by
Algebra, Lemma \ref{algebra-lemma-colimit-flat}.
Thus (3) $\Leftrightarrow$ (2).

\medskip\noindent
Assume (2). Suppose that $N \subset M$ with $M$ a flat $A$-module.
We will prove that $N$ is flat.
We can write $M = \colim M_i$ with each $M_i$ finite free, see
Algebra, Theorem \ref{algebra-theorem-lazard}.
Setting $N_i \subset M_i$ the inverse image of $N$ we see that
$N = \colim N_i$. By
Algebra, Lemma \ref{algebra-lemma-colimit-flat}.
it suffices to prove $N_i$ is flat and we reduce
to the case $M = R^{\oplus n}$. In this case
the module $N$ has a finite filtration by the submodules
$R^{\oplus j} \cap N$ whose subquotients are ideals.
By (2) these ideals are flat and hence $N$ is flat by
Algebra, Lemma \ref{algebra-lemma-flat-ses}. Thus (2) $\Rightarrow$ (4).

\medskip\noindent
Assume $A$ satisfies (1) and let $\mathfrak p \subset A$ be a
prime ideal. By
Lemmas \ref{lemma-when-weakly-etale} and \ref{lemma-weak-dimension-goes-up}
we see that $A_\mathfrak p$ satisfies (1). We will show $A$ is a valuation ring
if $A$ is a local ring satisfying (3). Let $f \in \mathfrak m$
be a nonzero element. Then $(f)$ is a flat nonzero module generated by
one element. Hence it is a free $A$-module by
Algebra, Lemma \ref{algebra-lemma-finite-flat-local}.
It follows that $f$ is a nonzerodivisor and $A$ is a domain.
If $I \subset A$ is a finitely generated ideal, then we similarly
see that $I$ is a finite free $A$-module, hence (by considering the
rank) free of rank $1$ and $I$ is a principal ideal. Thus $A$ is a
valuation ring by
Algebra, Lemma \ref{algebra-lemma-characterize-valuation-ring}.
Thus (1) $\Rightarrow$ (5).

\medskip\noindent
Assume (5). Let $I \subset A$ be a finitely generated ideal.
Then $I_\mathfrak p \subset A_\mathfrak p$ is a finitely generated ideal
in a valuation ring, hence principal
(Algebra, Lemma \ref{algebra-lemma-characterize-valuation-ring}), hence flat.
Thus $I$ is flat by
Algebra, Lemma \ref{algebra-lemma-flat-localization}.
Thus (5) $\Rightarrow$ (3). This finishes the proof of the lemma.
\end{proof}